\documentclass[twocolumn]{aastex631}

\usepackage{amsmath}
\usepackage{multirow}
\usepackage{natbib}
\usepackage{graphicx} 
\usepackage{aas_macros}

\begin{document}

In this paper, we have addressed the fundamental question of the quantum stability of Degenerate Higher-Order Scalar-Tensor (DHOST) theories. These theories are constructed to be free of the catastrophic Ostrogradsky ghost instability at the classical level through the imposition of specific algebraic degeneracy conditions. However, the robustness of this mechanism against quantum corrections has remained an open question. We proposed and executed a direct dynamical test of this stability by investigating the behavior of the degeneracy conditions under the Renormalization Group (RG) flow. We conceptualized the space of theories as a landscape where the degeneracy conditions define a "healthy" submanifold, and we tested whether this submanifold is an invariant of the quantum dynamics.

To perform this test, we computed the complete set of one-loop beta-functionals for the general Type Ia class of DHOST theories. Utilizing the path integral formalism with the background field method and heat kernel techniques, we determined how the defining functions of the theory evolve with the energy scale. Our central result is unambiguous: the RG flow vector is generically not tangent to the degeneracy manifold. We found that the quantum corrections generated for the functions that are classically constrained, such as $A_4$ and $A_5$, do not conspire to maintain the algebraic relations required for degeneracy. This demonstrates that quantum fluctuations systematically drive the theory away from the healthy, ghost-free subspace.

The physical implication of this finding is profound. The Ostrogradsky ghost, which is successfully projected out of the classical theory, is dynamically resurrected by quantum loop effects. A theory that begins precisely on the degeneracy manifold at a given energy scale will inevitably be pushed off it as the scale changes, causing the ghost degree of freedom to reappear with a non-zero kinetic term. The re-emergence of this pathological mode, with its associated negative energy, renders the quantum vacuum unstable and signals a fundamental inconsistency of the theory at the quantum level.

From our results, we conclude that the classical stability mechanism of Type Ia DHOST theories is an unprotected fine-tuning, not a fundamental feature protected by an underlying symmetry. The degeneracy conditions are only satisfied at a specific energy scale and must be continually readjusted order-by-order in perturbation theory, which is the hallmark of a non-renormalizable theory. Consequently, these theories cannot be considered as self-consistent, UV-complete quantum theories of gravity. They must instead be interpreted as low-energy effective field theories (EFTs), whose validity is limited by a cutoff scale. This scale is naturally associated with the energy at which the quantum-generated ghost mass becomes significant, signaling the breakdown of the EFT and the need for new physics. Our work thus establishes a powerful theoretical argument for the intrinsic quantum instability of this class of modified gravity theories, highlighting the crucial role that quantum consistency must play in the search for a viable theory of gravity beyond General Relativity.

\end{document}
                